\section{Our Approach: Prior-Guided Reinforcement Learning with Diffusion}
\label{sec:our_approach}

The optimal control formulation~\eqref{eq:oc_problem} provides a principled objective
for vaccine allocation on contact networks. However, direct solution via classical
optimal control methods is computationally intractable due to: (1) the exponentially
large state space ($6^N$ discrete states for $N$ individuals), (2) stochastic dynamics
requiring expectations over $2^N$ possible outcomes at each time step, and (3)
network-coupled transitions that violate the separability assumptions of standard
algorithms. While population-level ODE models with aggregated compartments can be
solved efficiently using direct collocation methods~\citep{师兄论文}, our individual-level
CTMP on networks with $N \sim 10^3$ nodes necessitates approximate solution methods.

We therefore adopt \textbf{reinforcement learning (RL)} to approximate the optimal
control solution. RL circumvents the computational barriers by: (1) using function
approximation (neural networks) to represent policies in high dimensions, (2) estimating
expectations via Monte Carlo sampling from simulated trajectories, and (3) optimizing
policies directly through gradient-based methods without solving Bellman equations
globally. However, naïve RL suffers from poor sample efficiency in this high-dimensional,
long-horizon problem ($N=2000$ individuals, $T=600$ time steps).

Our key insight is to leverage \emph{tractable} population-level optimal control
solutions as \textbf{priors} to guide individual-level RL. Specifically, we propose
a three-stage framework (Figure~\ref{fig:pipeline}):

\begin{enumerate}
    \item \textbf{Macro-level expert generation (M1):} Simulate group-level epidemic
    dynamics using deterministic ODE models and solve for optimal vaccination strategies
    using established methods. This produces expert trajectories $\{\tau_{\text{macro}}^{(k)}\}$
    at the population level.

    \item \textbf{Diffusion prior learning (M2):} Map macro-level actions to individual-level
    allocations using a degree-risk weighted lifting algorithm, then train a conditional
    diffusion model $p_\theta(a \mid s)$ to capture the distribution of expert strategies.

    \item \textbf{Prior-guided reinforcement learning (M3):} Train a Proximal Policy
    Optimization (PPO) agent on the individual-level CTMP environment, regularized by
    KL divergence to the diffusion prior:
    \begin{equation}
    \mathcal{L} = \mathcal{L}_{\text{PPO}}(s,a,A) + \beta(t) \cdot D_{\text{KL}}(\pi_\phi(\cdot|s) \,\|\, p_\theta(\cdot|s)),
    \label{eq:prior_guided_loss}
    \end{equation}
    where $\beta(t)$ decays over training to gradually shift from imitation to optimization.
\end{enumerate}

This hybrid approach combines the theoretical insights of optimal control (via ODE
experts) with the scalability of RL (via neural network policies), while using diffusion
models to bridge the macro-micro gap and improve sample efficiency. We now detail each
stage.


% ============================================
\subsection{Macro-Level ODE Expert Generation}
\label{sec:ode_expert}

We begin by generating expert vaccination strategies at the population level using
the well-established ODE framework for compartmental epidemic models. This provides
high-quality prior knowledge that would be expensive to discover through exploration
in the high-dimensional CTMP environment.

\subsubsection{Population-Level ODE Dynamics}

Let $S_g(t), E_g(t), I_g(t), R_g(t), V_g(t), D_g(t)$ denote the population densities
in each compartment for group $g \in \{1,2,3\}$. The dynamics follow a system of
ordinary differential equations:
\begin{align}
\frac{dS_g}{dt} &= -\lambda_g(t) S_g - u_g(t) S_g, \\
\frac{dE_g}{dt} &= \lambda_g(t) S_g + (1-\varepsilon_g)\lambda_g(t) V_g - \sigma_g E_g, \\
\frac{dI_g}{dt} &= \sigma_g E_g - \gamma_g I_g, \\
\frac{dR_g}{dt} &= \gamma_g(1-\mu_g) I_g, \\
\frac{dV_g}{dt} &= u_g(t) S_g - (1-\varepsilon_g)\lambda_g(t) V_g, \\
\frac{dD_g}{dt} &= \gamma_g \mu_g I_g,
\end{align}
where the population-level force of infection is given by the mean-field approximation:
\begin{equation}
\lambda_g(t) = \beta_g \sum_{h=1}^3 C_{gh} \frac{I_h(t)}{N_h},
\label{eq:ode_foi}
\end{equation}
with $C_{gh}$ denoting the contact matrix between groups and $N_h$ the size of group $h$.
The control $u_g(t) \in [0,1]$ represents the vaccination rate for group $g$, subject
to the aggregate constraint:
\begin{equation}
\sum_{g=1}^3 u_g(t) S_g(t) \leq V_{\max}.
\label{eq:ode_constraint}
\end{equation}

\subsubsection{Expert Strategy Generation}

We generate expert trajectories using three canonical vaccination strategies, which
represent different policy priorities established in prior work~\citep{师兄论文}:

\begin{itemize}
    \item \textbf{High-risk priority:} Allocate vaccines sequentially to groups
    $g=2$ (high-risk), then $g=3$ (high-contact), then $g=1$ (baseline). This
    strategy prioritizes mortality reduction by protecting vulnerable populations.

    \item \textbf{High-contact priority:} Allocate to groups $g=3, 2, 1$ in order.
    This strategy prioritizes transmission interruption by vaccinating those with
    highest contact rates.

    \item \textbf{Uniform allocation:} Distribute vaccines proportionally across
    all groups according to their susceptible populations.
\end{itemize}

For each strategy and each random initial condition (varying initial infections,
network realizations), we simulate the ODE system forward for $T=600$ days. Each
simulation yields a macro-level trajectory:
\begin{equation}
\tau_{\text{macro}}^{(k)} = \left\{
(S_g^{(k)}(t), E_g^{(k)}(t), I_g^{(k)}(t), R_g^{(k)}(t), V_g^{(k)}(t), D_g^{(k)}(t),
u_g^{(k)}(t))_{g=1,2,3}
\right\}_{t=0}^{T-1}.
\end{equation}

We generate $K=1000$ such expert trajectories (distributed across the three strategies
and diverse initial conditions) using a modified version of the ProtectorPrevent
simulator~\citep{protectorprevent}. These trajectories serve as demonstrations of
effective vaccination policies at the population level, which we will subsequently
map to the individual level.


% ============================================
\subsection{Macro-to-Micro Lifting Algorithm}
\label{sec:lifting}

The macro-level expert trajectories $\{\tau_{\text{macro}}^{(k)}\}$ prescribe
group-level vaccination rates $u_g(t)$ but do not specify \emph{which individuals}
within each group should be vaccinated. To transfer this knowledge to the individual-level
CTMP, we require a \textbf{lifting algorithm} that maps group allocations to
individual probabilities while respecting network structure and feasibility constraints.

\subsubsection{Degree-Risk Weighted Lifting}

Given a macro-level control $U_t = (u_1(t), u_2(t), u_3(t))$ and the current
individual-level state $s_t = (X_1(t), \ldots, X_N(t))$, we compute individual
vaccination probabilities $a_t = (a_1(t), \ldots, a_N(t))$ as follows.

\paragraph{Step 1: Compute individual weights.}
For each susceptible individual $i$ with $X_i(t) = S$ and group label $g_i$, define
a priority weight combining network centrality and risk:
\begin{equation}
w_i = \alpha \cdot \frac{k_i}{\max_j k_j} + (1-\alpha) \cdot \mathbb{I}\{g_i = 2\},
\label{eq:weight}
\end{equation}
where $k_i = |\mathcal{N}(i)|$ is the degree of individual $i$, and $\mathbb{I}\{g_i=2\}$
is an indicator for the high-risk group. We use $\alpha = 0.6$ to balance network
effects (degree) with demographic risk. This weighting reflects the principle that
high-degree individuals act as ``super-spreaders'' in the network, while high-risk
individuals face greater morbidity and mortality.

\paragraph{Step 2: Within-group normalization.}
For each group $g$, let $V_g = \{i : X_i(t)=S, g_i=g\}$ denote the set of susceptible
individuals in group $g$. Normalize the weights within each group:
\begin{equation}
\tilde{w}_i = \frac{w_i}{\sum_{j \in V_g} w_j}, \quad \forall i \in V_g.
\end{equation}

\paragraph{Step 3: Allocate according to macro control.}
The desired vaccination budget for group $g$ is:
\begin{equation}
B_g = u_g(t) \cdot |V_g|.
\end{equation}
We distribute this budget among individuals in $V_g$ proportionally to their normalized
weights:
\begin{equation}
a_i^{\text{raw}}(t) = \tilde{w}_i \cdot B_g, \quad \forall i \in V_g.
\label{eq:raw_action}
\end{equation}

\paragraph{Step 4: Feasibility projection.}
The raw actions $a_i^{\text{raw}}$ may violate the global supply constraint
$\sum_i a_i \leq V_{\max}$. We project onto the feasible set:
\begin{equation}
a_i(t) = \min\left(a_i^{\text{raw}}(t), 1\right) \cdot \rho,
\quad \text{where} \quad
\rho = \min\left(1, \frac{V_{\max}}{\sum_i a_i^{\text{raw}}(t)}\right).
\label{eq:projection}
\end{equation}
This ensures $a_i(t) \in [0,1]$ and $\sum_i a_i(t) \leq V_{\max}$.

The complete lifting procedure is summarized in Algorithm~\ref{alg:lifting}.

\begin{algorithm}[t]
\caption{Macro-to-Micro Lifting Algorithm}
\label{alg:lifting}
\begin{algorithmic}[1]
\STATE \textbf{Input:} Macro control $U_t = (u_1, u_2, u_3)$, state $s_t$, graph $\mathcal{G}$, $V_{\max}$
\STATE \textbf{Output:} Individual actions $a_t = (a_1, \ldots, a_N)$
\STATE
\STATE $a_i \gets 0$ for all $i \in V$ \hfill $\triangleright$ Initialize
\FOR{each group $g \in \{1,2,3\}$}
    \STATE $V_g \gets \{i : X_i(t)=S, g_i=g\}$ \hfill $\triangleright$ Susceptible in group $g$
    \IF{$|V_g| = 0$}
        \STATE \textbf{continue}
    \ENDIF
    \STATE
    \STATE \textcolor{blue}{\# Step 1-2: Compute normalized weights}
    \FOR{$i \in V_g$}
        \STATE $w_i \gets 0.6 \cdot k_i / \max_j k_j + 0.4 \cdot \mathbb{I}\{g_i=2\}$
    \ENDFOR
    \STATE $\tilde{w}_i \gets w_i / \sum_{j \in V_g} w_j$ for all $i \in V_g$
    \STATE
    \STATE \textcolor{blue}{\# Step 3: Allocate budget}
    \STATE $B_g \gets u_g \cdot |V_g|$
    \FOR{$i \in V_g$}
        \STATE $a_i^{\text{raw}} \gets \tilde{w}_i \cdot B_g$
    \ENDFOR
\ENDFOR
\STATE
\STATE \textcolor{blue}{\# Step 4: Feasibility projection}
\STATE $\rho \gets \min(1, V_{\max} / \sum_i a_i^{\text{raw}})$
\FOR{$i \in V$}
    \STATE $a_i \gets \min(a_i^{\text{raw}}, 1) \cdot \rho$
\ENDFOR
\RETURN $a_t = (a_1, \ldots, a_N)$
\end{algorithmic}
\end{algorithm}

\subsubsection{Lifted Expert Dataset}

Applying Algorithm~\ref{alg:lifting} to each macro-level expert trajectory
$\tau_{\text{macro}}^{(k)}$ yields a corresponding individual-level trajectory:
\begin{equation}
\tau_{\text{micro}}^{(k)} = \left\{
(s_t^{(k)}, a_t^{(k)})
\right\}_{t=0}^{T-1},
\end{equation}
where $s_t^{(k)}$ is obtained by simulating the CTMP (Algorithm~1 in Section~3) under
the lifted actions $a_t^{(k)}$. The collection
$\mathcal{D} = \{\tau_{\text{micro}}^{(k)}\}_{k=1}^K$ forms a dataset of individual-level
expert demonstrations, which we use to train the diffusion prior in the next stage.


% ============================================
\subsection{Conditional Diffusion Model for Strategy Prior}
\label{sec:diffusion}

Rather than using the lifted expert data directly for behavior cloning (which would
yield a deterministic policy), we train a \textbf{conditional diffusion model} to
capture the \emph{distribution} of expert strategies. This probabilistic representation
has two advantages: (1) it provides a smooth prior for KL regularization in RL, and
(2) it enables sampling diverse strategies that generalize beyond the expert demonstrations.

\subsubsection{Denoising Diffusion Probabilistic Model}

We adopt the Denoising Diffusion Probabilistic Model (DDPM) framework~\citep{ho2020denoising}.
Let $a \in \mathbb{R}^N$ denote an individual-level action vector (vaccination probabilities).
The forward diffusion process gradually adds Gaussian noise over $T_{\text{diff}}$ steps:
\begin{equation}
q(a^{(1:T_{\text{diff}})} \mid a^{(0)}) = \prod_{t=1}^{T_{\text{diff}}} q(a^{(t)} \mid a^{(t-1)}),
\quad
q(a^{(t)} \mid a^{(t-1)}) = \mathcal{N}(a^{(t)}; \sqrt{1-\beta_t} a^{(t-1)}, \beta_t I),
\end{equation}
where $\{\beta_t\}$ is a variance schedule (we use $T_{\text{diff}}=1000$ and linear
schedule $\beta_t \in [10^{-4}, 0.02]$). The reverse process learns to denoise:
\begin{equation}
p_\theta(a^{(0:T_{\text{diff}})}) = p(a^{(T_{\text{diff}})}) \prod_{t=1}^{T_{\text{diff}}}
p_\theta(a^{(t-1)} \mid a^{(t)}),
\end{equation}
where $p(a^{(T_{\text{diff}})}) = \mathcal{N}(0, I)$ and the denoising distribution
is parameterized as:
\begin{equation}
p_\theta(a^{(t-1)} \mid a^{(t)}) = \mathcal{N}(a^{(t-1)}; \mu_\theta(a^{(t)}, t), \Sigma_\theta(a^{(t)}, t)).
\end{equation}

Following~\cite{ho2020denoising}, we use a simplified objective that predicts the
noise $\epsilon \sim \mathcal{N}(0,I)$ added at each step:
\begin{equation}
\mathcal{L}_{\text{diff}} = \mathbb{E}_{a^{(0)} \sim \mathcal{D}, t, \epsilon}
\left[ \|\epsilon - \epsilon_\theta(a^{(t)}, t, s)\|^2 \right],
\label{eq:ddpm_loss}
\end{equation}
where $a^{(t)} = \sqrt{\bar{\alpha}_t} a^{(0)} + \sqrt{1-\bar{\alpha}_t} \epsilon$
with $\bar{\alpha}_t = \prod_{i=1}^t (1-\beta_i)$, and $s$ is the epidemic state
used for conditioning.

\subsubsection{Conditional Architecture}

To condition the diffusion model on the epidemic state $s_t$, we extract a
$d_{\text{feat}}$-dimensional state feature vector:
\begin{equation}
\phi(s_t) = \left(
\frac{|S_g|}{N_g}, \frac{|E_g|}{N_g}, \frac{|I_g|}{N_g}, \frac{|R_g|}{N_g},
\frac{|V_g|}{N_g}, \frac{|D_g|}{N_g}
\right)_{g=1,2,3} \in \mathbb{R}^{18},
\end{equation}
where $|S_g| = |\{i : X_i(t)=S, g_i=g\}|$ counts individuals in compartment $S$ and
group $g$, and similarly for other compartments. We concatenate $\phi(s_t)$ with
additional features (time step $t/T$, total supply $V_{\max}$) to obtain
$c_t \in \mathbb{R}^{20}$.

The denoising network $\epsilon_\theta(a^{(t)}, t, c)$ is implemented as a
Transformer~\citep{vaswani2017attention} with the following components:
\begin{itemize}
    \item \textbf{Input embedding:} Project the noisy action $a^{(t)} \in \mathbb{R}^N$
    to dimension $d_{\text{model}}=128$ via a linear layer.

    \item \textbf{Conditioning:} Inject the state features $c$ and diffusion timestep $t$
    via sinusoidal positional encoding and addition to the input embedding.

    \item \textbf{Transformer layers:} 4 layers of multi-head self-attention (4 heads)
    with feedforward dimension 512 and LayerNorm.

    \item \textbf{Output projection:} Map the final hidden states to $\mathbb{R}^N$
    to predict the noise $\epsilon$.
\end{itemize}

\subsubsection{Training Procedure}

We train the diffusion model on the lifted expert dataset $\mathcal{D} =
\{(s_t^{(k)}, a_t^{(k)})\}$ by minimizing the DDPM loss~\eqref{eq:ddpm_loss}. Training
uses the AdamW optimizer with learning rate $3 \times 10^{-4}$, batch size 128, and
cosine annealing schedule over 50 epochs. After training, we obtain a conditional
generative model $p_\theta(a \mid s)$ that can sample vaccination strategies conditioned
on arbitrary epidemic states.

To use this model as a prior during RL, we evaluate the log-probability
$\log p_\theta(a \mid s)$ via importance sampling~\citep{song2021scorebased}, which
allows computing KL divergence for regularization (see Section~\ref{sec:rl}).


% ============================================
\subsection{Prior-Guided Reinforcement Learning}
\label{sec:rl}

Having constructed a diffusion-based prior $p_\theta(a \mid s)$ that captures expert
knowledge, we now train a policy $\pi_\phi(a \mid s)$ to optimize the individual-level
CTMP objective~\eqref{eq:oc_problem} while regularized by this prior. We employ
Proximal Policy Optimization (PPO)~\citep{schulman2017ppo} with KL regularization.

\subsubsection{Policy and Value Networks}

\paragraph{Policy network $\pi_\phi(a \mid s)$.}
The policy is parameterized as a Gaussian distribution with state-dependent mean and
fixed standard deviation:
\begin{equation}
\pi_\phi(a \mid s) = \mathcal{N}(a; \mu_\phi(s), \sigma^2 I),
\end{equation}
where $\mu_\phi : \mathcal{S} \to \mathbb{R}^N$ is a neural network. The network
architecture consists of:
\begin{itemize}
    \item State encoder: Maps the individual states and graph structure to a
    $d_{\text{hidden}}=256$-dimensional representation using a 3-layer MLP with
    group-level aggregation.
    \item Output head: Projects to $\mathbb{R}^N$ and applies sigmoid activation
    to ensure $\mu_\phi(s) \in [0,1]^N$.
\end{itemize}
We use $\sigma = 0.1$ as a fixed exploration noise level.

\paragraph{Value network $V_\psi(s)$.}
The baseline for variance reduction in PPO is a learned state-value function
$V_\psi : \mathcal{S} \to \mathbb{R}$, implemented as a 3-layer MLP with the same
state encoder as the policy network.

\subsubsection{PPO with KL Regularization}

Standard PPO maximizes the clipped surrogate objective:
\begin{equation}
\mathcal{L}_{\text{PPO}}(\phi) = \mathbb{E}_{\tau \sim \pi_{\phi_{\text{old}}}}
\left[
\min\left(
r_t(\phi) \hat{A}_t, \,
\text{clip}(r_t(\phi), 1-\epsilon, 1+\epsilon) \hat{A}_t
\right)
\right],
\label{eq:ppo_loss}
\end{equation}
where $r_t(\phi) = \pi_\phi(a_t \mid s_t) / \pi_{\phi_{\text{old}}}(a_t \mid s_t)$
is the importance sampling ratio, $\hat{A}_t$ is the generalized advantage estimate
(GAE)~\citep{schulman2015gae}, and $\epsilon=0.2$ is the clipping parameter.

We augment this with KL regularization to the diffusion prior:
\begin{equation}
\mathcal{L}_{\text{total}}(\phi) = \mathcal{L}_{\text{PPO}}(\phi)
+ \beta(n) \cdot \mathbb{E}_{s \sim \mathcal{D}_{\text{buffer}}}
\left[ D_{\text{KL}}(\pi_\phi(\cdot \mid s) \,\|\, p_\theta(\cdot \mid s)) \right]
- \lambda_{\text{ent}} \mathcal{H}(\pi_\phi),
\label{eq:total_loss}
\end{equation}
where:
\begin{itemize}
    \item $D_{\text{KL}}(\pi_\phi \| p_\theta)$ is the KL divergence from the RL
    policy to the diffusion prior, encouraging the policy to stay close to expert
    demonstrations early in training.

    \item $\beta(n) = \beta_0 \cdot \gamma^{n}$ is a decaying coefficient
    ($\beta_0=1.0$, $\gamma=0.995$ per epoch $n$), which gradually reduces prior
    influence as the policy improves through exploration.

    \item $\mathcal{H}(\pi_\phi)$ is the policy entropy for exploration, weighted
    by $\lambda_{\text{ent}}=0.01$.
\end{itemize}

The value function is trained by minimizing the mean-squared Bellman error:
\begin{equation}
\mathcal{L}_{\text{value}}(\psi) = \mathbb{E}_{(s_t,R_t) \sim \mathcal{D}_{\text{buffer}}}
\left[ (V_\psi(s_t) - R_t)^2 \right],
\end{equation}
where $R_t = \sum_{k=0}^{T-t-1} \gamma^k r_{t+k}$ is the discounted return
($\gamma=0.99$).

\subsubsection{Training Algorithm}

The complete training procedure is outlined in Algorithm~\ref{alg:ppo_training}.
At each iteration, we collect trajectories by rolling out the current policy
$\pi_\phi$ in the CTMP environment (Algorithm~1), compute advantages using GAE,
and update the policy and value networks via gradient descent on~\eqref{eq:total_loss}
and~\eqref{eq:value_loss}.

\begin{algorithm}[t]
\caption{Prior-Guided PPO Training}
\label{alg:ppo_training}
\begin{algorithmic}[1]
\STATE \textbf{Input:} Diffusion prior $p_\theta$, CTMP environment, $N_{\text{iter}}$ iterations
\STATE \textbf{Output:} Trained policy $\pi_\phi^*$
\STATE
\STATE Initialize policy $\pi_\phi$, value function $V_\psi$ with random weights
\FOR{iteration $n = 1, \ldots, N_{\text{iter}}$}
    \STATE \textcolor{blue}{\# Collect rollouts}
    \STATE $\mathcal{D}_{\text{buffer}} \gets \emptyset$
    \FOR{episode $e = 1, \ldots, N_{\text{episodes}}$}
        \STATE $s_0 \sim$ initial state distribution
        \FOR{$t = 0, \ldots, T-1$}
            \STATE $a_t \sim \pi_\phi(\cdot \mid s_t)$
            \STATE $s_{t+1} \sim P(\cdot \mid s_t, a_t)$ \hfill $\triangleright$ CTMP simulation
            \STATE $r_t \gets -(|I_t| + 10|D_t|)$
            \STATE $\mathcal{D}_{\text{buffer}} \gets \mathcal{D}_{\text{buffer}} \cup \{(s_t, a_t, r_t)\}$
        \ENDFOR
    \ENDFOR
    \STATE
    \STATE \textcolor{blue}{\# Compute advantages}
    \STATE Compute GAE advantages $\{\hat{A}_t\}$ using $V_\psi$
    \STATE
    \STATE \textcolor{blue}{\# Update policy with KL regularization}
    \STATE $\beta \gets \beta_0 \cdot 0.995^n$
    \FOR{PPO epoch $k = 1, \ldots, K_{\text{ppo}}$}
        \STATE Sample mini-batch from $\mathcal{D}_{\text{buffer}}$
        \STATE Update $\phi$ by gradient ascent on $\mathcal{L}_{\text{total}}(\phi)$ (Eq.~\ref{eq:total_loss})
        \STATE Update $\psi$ by gradient descent on $\mathcal{L}_{\text{value}}(\psi)$
    \ENDFOR
\ENDFOR
\RETURN $\pi_\phi$
\end{algorithmic}
\end{algorithm}

\subsubsection{Comparison Baselines}

To validate the effectiveness of the diffusion prior, we train three variants:
\begin{enumerate}
    \item \textbf{RL-only:} Standard PPO without prior ($\beta=0$ in Eq.~\ref{eq:total_loss}).
    \item \textbf{BC+RL:} Behavior cloning warmstart (supervised learning on $\mathcal{D}$)
    followed by PPO fine-tuning.
    \item \textbf{Diffusion+RL (Ours):} Prior-guided PPO with diffusion model and
    decaying KL regularization.
\end{enumerate}
Experimental comparisons are presented in Section~\ref{sec:experiments}.


% ============================================
\subsection{Implementation Details}
\label{sec:implementation}

All models are implemented in PyTorch. The CTMP simulator uses NumPy for efficient
state updates and NetworkX for graph operations. Training is conducted on a single
NVIDIA GPU (RTX 3090) with the following hyperparameters:

\begin{itemize}
    \item \textbf{Diffusion model:} $T_{\text{diff}}=1000$, $d_{\text{model}}=128$,
    4 Transformer layers, learning rate $3 \times 10^{-4}$, 50 epochs.
    \item \textbf{PPO:} Learning rate $3 \times 10^{-4}$, $N_{\text{episodes}}=10$
    per iteration, $K_{\text{ppo}}=10$ optimization epochs, GAE $\lambda=0.95$,
    discount $\gamma=0.99$, clip $\epsilon=0.2$.
    \item \textbf{KL regularization:} $\beta_0=1.0$, decay $\gamma=0.995$ per iteration.
\end{itemize}

Training RL-only requires $\sim 200$ iterations ($\approx 2000$ episodes) to converge,
taking approximately 12 hours. Diffusion+RL converges in $\sim 100$ iterations
($\approx 1000$ episodes, 6 hours), demonstrating improved sample efficiency.
